\documentclass[a4paper,11pt,english]{article}

\usepackage[margin=2.5cm]{geometry}
\usepackage[utf8]{inputenc}
\usepackage[T1]{fontenc}
\usepackage[english]{babel}
\usepackage{amsmath}
\usepackage{amssymb}
\usepackage{amsthm}
\usepackage[colorlinks,urlcolor=blue,citecolor=blue,linkcolor=blue]{hyperref}

\newtheorem{theorem}{Theorem}[section]
\newtheorem{lemma}[theorem]{Lemma}
\newtheorem{corollary}[theorem]{Corollary}

\newcommand{\E}{\mathbb E}
\newcommand{\Prob}{\mathbb P}
\newcommand{\Var}{\operatorname{Var}}
\newcommand{\Cov}{\operatorname{Cov}}
\newcommand{\diag}{\operatorname{diag}}
\newcommand{\norm}[1]{\left\lVert #1\right\rVert}
\newcommand{\R}{\mathbb R}

\setlength{\emergencystretch}{3em}

\title{Matched Gaussian Chaos and Contrast Approximation}
\author{Th\'eo Voldoire and Nic Fishman}
\date{}

\begin{document}

\maketitle

\begin{abstract}
Searching over $s$ controls reduces the largest regression coefficient to
a covariance-normalized bilinear functional of two empirical score
vectors.  Under relative sub-Gaussianity and
$[s\log(ep/s)]^4=o(n)$, a matched Gaussian law approximates this
functional while retaining its fourth-moment dependence.  A fixed
two-dimensional rotation writes every support value as one projection
energy minus another.  Any support chosen from the first contrast then
approximates the search up to its sparse-projection gap and one
$s$-dimensional remainder energy.  Conditional covariance and mean bounds
make the remainder negligible for arbitrary control covariance, so a
support with a negligible projection gap gives a stochastic
approximation.  Under Kronecker Gaussian score covariance, the remainder
has an exact conditional noncentral chi-square law.  A bounded restricted
condition number gives a sorting-based
constant-factor approximation and a computable forward-selection
certificate, while local isotropy recovers the top-$s$ product formula.
\end{abstract}

\section{Regression search}

Adding controls changes the coefficient on $x_0$ through two score
projections.  The reduction below isolates this dependence in a bilinear
search over supports, whose covariance geometry is the only obstacle to
the coordinatewise top-$s$ construction.

For independent and identically distributed observations
$i=1,\ldots,n$, consider the centered linear model
\begin{equation}\label{eq:model}
        y^{(i)}
        =\beta_0x_0^{(i)}
        +(X_{-0}^{(i)})^\top\beta_{-0}
        +\varepsilon^{(i)}.
\end{equation}
The same letters without row superscripts denote the stacked sample
vectors, and $X_{-0}$ denotes the matrix with rows $X_{-0}^{(i)}$.
The pair $(x_0^{(i)},X_{-0}^{(i)})$ is jointly sub-Gaussian relative to
its covariance, and $\varepsilon^{(i)}$ is centered, sub-Gaussian, and
independent of the design.  Write
\begin{equation*}
        \Var(x_0^{(i)})=\sigma_0^2,
        \qquad
        \Cov(X_{-0}^{(i)})=C,
        \qquad
        d=\Cov(X_{-0}^{(i)},x_0^{(i)}),
\end{equation*}
where $\sigma_0>0$.  Every standardized linear form of the design and
noise has uniformly bounded $\psi_2$ norm.  Rescaling a control leaves
every selected column span unchanged, so we assume throughout that
$C_{jj}=1$ for every $j$.

Write $[p]=\{1,\ldots,p\}$.  For $S\subseteq[p]$, let $M_S$ project
orthogonally off the column span of $X_{-0,S}$, and define
\begin{equation}\label{eq:searched-coefficient}
        \widehat\beta_{0,S}
        =\frac{x_0^\top M_Sy}{x_0^\top M_Sx_0}.
\end{equation}
The target is the supremum of \eqref{eq:searched-coefficient} over
$|S|=s$.  Define the residual outcome and its score mean by
\begin{equation}\label{eq:residual-outcome}
\begin{aligned}
        r^{(i)}
        &=y^{(i)}-\beta_0x_0^{(i)}
        -\frac{d^\top\beta_{-0}}{\sigma_0^2}x_0^{(i)},\\
        h&=\Cov(X_{-0}^{(i)},r^{(i)})
        =C\beta_{-0}
        -\frac{d^\top\beta_{-0}}{\sigma_0^2}d,
        &\sigma_r^2&=\Var(r^{(i)}).
\end{aligned}
\end{equation}
Then $\Cov(x_0^{(i)},r^{(i)})=0$.  We assume $\sigma_r>0$.

For $v\in\R^p$ and $1\le k\le2s$, define the
covariance-normalized sparse norm
\begin{equation}\label{eq:sparse-norm}
        \norm{v}_{C,k}
        =\sup_{J\subseteq[p],\,|J|\le k}
        \left(v_J^\top C_{JJ}^{-1}v_J\right)^{1/2},
\end{equation}
and assume that the blocks $C_{JJ}$ are positive definite for
$|J|\le2s$.  Write
\begin{equation}\label{eq:scales}
        L=\log(ep/s),
        \qquad
        \delta_{s,n,p}=\sqrt{sL/n}.
\end{equation}
The population score means are diffuse when
\begin{equation}\label{eq:proxy-energy}
        \norm{d}_{C,2s}
        \lesssim\sigma_0\delta_{s,n,p},
        \qquad
        \norm{h}_{C,2s}
        \lesssim\sigma_r\delta_{s,n,p}.
\end{equation}

For $a,b\in\R^p$ and $S\subseteq[p]$ with $|S|\le2s$, put
\begin{equation}\label{eq:search-functional}
        Q_S(a,b)=-a_S^\top C_{SS}^{-1}b_S,
        \qquad
        T_C(a,b)=\max_{|S|=s}Q_S(a,b).
\end{equation}
For centered jointly Gaussian $a$ and $b$, each
$Q_S(a,b)-\E Q_S(a,b)$ is a second-order Gaussian chaos, and $T_C(a,b)$
is the maximum of the corresponding uncentered quadratic forms.

\section{Contrast approximation}

The bilinear search has a one-contrast reduction.  For positive
$\sigma_a,\sigma_b$ and $a,b\in\R^p$, define the search contrast and its
remainder by
\begin{equation}\label{eq:contrasts}
        G=\frac{1}{\sqrt2}\left(\frac a{\sigma_a}
        -\frac b{\sigma_b}\right),
        \qquad
        H=\frac{1}{\sqrt2}\left(\frac a{\sigma_a}
        +\frac b{\sigma_b}\right).
\end{equation}
Polarization gives, for every support $S$,
\begin{equation}\label{eq:contrast-identity}
        Q_S(a,b)
        =\frac{\sigma_a\sigma_b}{2}
        \left(
        G_S^\top C_{SS}^{-1}G_S
        -H_S^\top C_{SS}^{-1}H_S
        \right).
\end{equation}
Thus a support chosen from $G$ leaves only one remainder energy to
control, rather than a second maximum over the support family.

\begin{theorem}[Contrast approximation]
\label{thm:contrast-approximation}
For every $a,b\in\R^p$ and every size-$s$ support $\widehat S$,
\begin{equation}\label{eq:contrast-transfer}
\begin{split}
        0\le T_C(a,b)-Q_{\widehat S}(a,b)
        \le\frac{\sigma_a\sigma_b}{2}\bigg\{
        &\norm{G}_{C,s}^2
        -G_{\widehat S}^\top C_{\widehat S\widehat S}^{-1}
        G_{\widehat S}\\
        &+H_{\widehat S}^\top C_{\widehat S\widehat S}^{-1}
        H_{\widehat S}\bigg\}.
\end{split}
\end{equation}
In particular, let $S_G$ maximize
$G_S^\top C_{SS}^{-1}G_S$ over $|S|=s$, with ties resolved
deterministically.  Then
\begin{equation}\label{eq:contrast-oracle-bound}
\begin{split}
        0
        &\le T_C(a,b)-Q_{S_G}(a,b)\\
        &\le
        \frac{\sigma_a\sigma_b}{2}
        H_{S_G}^\top C_{S_GS_G}^{-1}H_{S_G},
\end{split}
\end{equation}
and this remainder also bounds
$\sigma_a\sigma_b\norm{G}_{C,s}^2/2-T_C(a,b)$ from above.
\end{theorem}

\begin{proof}
The remainder energy in \eqref{eq:contrast-identity} is nonnegative, so
\begin{equation*}
        T_C(a,b)
        \le\frac{\sigma_a\sigma_b}{2}\norm{G}_{C,s}^2.
\end{equation*}
Subtracting \eqref{eq:contrast-identity} at $\widehat S$ proves
\eqref{eq:contrast-transfer}.  Taking $\widehat S=S_G$ removes the
optimization gap, while evaluating the maximum at $S_G$ proves the final
upper bound.
\end{proof}

The stochastic result uses the following sparse conditional covariance
domination: for a constant $K$,
\begin{equation}\label{eq:conditional-passive-covariance}
        \Cov(H_J\mid G)\preceq K C_{JJ}
        \qquad\text{for every }|J|\le s.
\end{equation}
Although \eqref{eq:conditional-passive-covariance} is stated over the
support family, the proof uses it only at the support selected from $G$.

\begin{corollary}[Stochastic contrast approximation]
\label{cor:stochastic-contrast}
Let $a,b$ be random vectors, let $\widehat S$ be any $G$-measurable
size-$s$ support, and suppose
\eqref{eq:conditional-passive-covariance} holds with $K=O(1)$.  Then
\begin{equation}\label{eq:passive-contrast-size}
        H_{\widehat S}^\top C_{\widehat S\widehat S}^{-1}H_{\widehat S}
        =O_{\Prob}\left(
        s+\norm{\E[H\mid G]}_{C,s}^2
        \right).
\end{equation}
If $L\to\infty$,
\begin{equation}\label{eq:passive-mean-condition}
        \norm{\E[H\mid G]}_{C,s}^2=o_{\Prob}(sL),
\end{equation}
and
\begin{equation}\label{eq:generic-projection-gap}
        \norm{G}_{C,s}^2
        -G_{\widehat S}^\top C_{\widehat S\widehat S}^{-1}
        G_{\widehat S}=o_{\Prob}(sL),
\end{equation}
then
\begin{equation}\label{eq:stochastic-contrast-computed}
        T_C(a,b)
        =Q_{\widehat S}(a,b)
        +o_{\Prob}(\sigma_a\sigma_b sL).
\end{equation}
For $\widehat S=S_G$, the optimization gap vanishes and
\begin{equation}\label{eq:stochastic-contrast-oracle}
        T_C(a,b)
        =Q_{S_G}(a,b)+o_{\Prob}(\sigma_a\sigma_b sL)
        =\frac{\sigma_a\sigma_b}{2}\norm{G}_{C,s}^2
        +o_{\Prob}(\sigma_a\sigma_b sL).
\end{equation}
\end{corollary}

\begin{proof}
Conditioning on $G$ fixes $\widehat S$, whose conditional expected
remainder energy is at most
\begin{equation*}
        Ks+
        \E[H_{\widehat S}\mid G]^\top
        C_{\widehat S\widehat S}^{-1}
        \E[H_{\widehat S}\mid G].
\end{equation*}
Conditional Markov inequality proves \eqref{eq:passive-contrast-size},
and Theorem~\ref{thm:contrast-approximation} proves the remaining claims.
\end{proof}

For a Gaussian pair with Kronecker covariance, the remainder has an exact
law.  Suppose, for $-1<\gamma<1$, that
\begin{equation}\label{eq:separable-gaussian-channels}
        \begin{pmatrix}a/\sigma_a\\ b/\sigma_b\end{pmatrix}
        \sim N\!\left(
        \begin{pmatrix}m_a\\m_b\end{pmatrix},
        \begin{pmatrix}1&\gamma\\\gamma&1\end{pmatrix}\otimes C
        \right).
\end{equation}
The contrasts $G$ and $H$ are independent Gaussian vectors with
covariances $(1-\gamma)C$ and $(1+\gamma)C$, respectively.

\begin{corollary}[Gaussian contrasts]
\label{cor:gaussian-contrasts}
Under \eqref{eq:separable-gaussian-channels}, every $G$-measurable
size-$s$ support $\widehat S$ satisfies, conditionally on $G$,
\begin{equation}\label{eq:passive-chi-square}
        \frac{H_{\widehat S}^\top
        C_{\widehat S\widehat S}^{-1}H_{\widehat S}}
        {1+\gamma}
        \sim\chi_s^2\!\left(
        \frac{(\E H)_{\widehat S}^\top
        C_{\widehat S\widehat S}^{-1}(\E H)_{\widehat S}}
        {1+\gamma}
        \right).
\end{equation}
In the centered case, for every positive definite $C$,
\begin{equation}\label{eq:centered-gaussian-contrast}
\begin{split}
        T_C(a,b)
        &=Q_{S_G}(a,b)+O_{\Prob}(\sigma_a\sigma_b s)\\
        &=\frac{\sigma_a\sigma_b}{2}\norm{G}_{C,s}^2
        +O_{\Prob}(\sigma_a\sigma_b s).
\end{split}
\end{equation}
\end{corollary}

\begin{proof}
Independence makes the conditional law of $H$ equal to its marginal law.
Whitening its selected $s$-block gives
\eqref{eq:passive-chi-square}; Theorem~\ref{thm:contrast-approximation}
then gives \eqref{eq:centered-gaussian-contrast}.
\end{proof}

\section{Gaussian score approximation}

Define the centered row score and its covariance by
\begin{equation}\label{eq:matched-score}
        U^{(i)}=
        \begin{pmatrix}
        x_0^{(i)}X_{-0}^{(i)}-d\\
        r^{(i)}X_{-0}^{(i)}-h
        \end{pmatrix},
        \qquad
        \Sigma_U=\Var(U^{(i)}),
\end{equation}
and collect the empirical scores in
\begin{equation}\label{eq:score-sum}
        Z_n=\begin{pmatrix}Z_{1n}\\Z_{2n}\end{pmatrix}
        :=\begin{pmatrix}
        X_{-0}^\top x_0/\sqrt n\\
        X_{-0}^\top r/\sqrt n
        \end{pmatrix}
        =\begin{pmatrix}\sqrt n\,d\\\sqrt n\,h\end{pmatrix}
        +\frac1{\sqrt n}\sum_{i=1}^nU^{(i)}.
\end{equation}
The matched Gaussian score vector is
\begin{equation}\label{eq:matched-gaussian}
        Z=\begin{pmatrix}Z_1\\Z_2\end{pmatrix}
        \sim N\!\left(
        \begin{pmatrix}\sqrt n\,d\\\sqrt n\,h\end{pmatrix},
        \Sigma_U\right).
\end{equation}
The covariance $\Sigma_U$ contains the fourth-moment contribution of the
design.

For a pair $z=(z_1,z_2)$, write
\begin{equation}\label{eq:standardized-sparse-norm}
        \norm{z}_s
        =\frac{\norm{z_1}_{C,s}}{\sigma_0}
        +\frac{\norm{z_2}_{C,s}}{\sigma_r}.
\end{equation}

For real random variables $\zeta,\widetilde\zeta$, write
\begin{equation*}
        d_3(\zeta,\widetilde\zeta)
        =\sup_{\varphi}
        |\E\varphi(\zeta)-\E\varphi(\widetilde\zeta)|.
\end{equation*}
The supremum is over all three times continuously differentiable
$\varphi$ such that
$\max_{0\le j\le3}\norm{\varphi^{(j)}}_\infty\le1$.

\begin{lemma}[Sparse score moments]
\label{lem:sparse-score-moments}
Let $\widetilde U^{(1)},\ldots,\widetilde U^{(n)}$ be independent, with
each $\widetilde U^{(i)}$ distributed either as $U^{(i)}$ or as a
centered Gaussian vector with covariance $\Sigma_U$.  Suppose
$sL=o(n)$.  Uniformly over all such choices and every integer
$1\le q\le6$,
\begin{equation}\label{eq:sparse-score-moments}
\begin{split}
        \E\norm{\widetilde U^{(i)}}_s^q&\lesssim(sL)^{q/2},\\
        \E\left|
        \begin{pmatrix}\sqrt n\,d\\\sqrt n\,h\end{pmatrix}
        +\frac1{\sqrt n}\sum_{i=1}^n\widetilde U^{(i)}
        \right|_s^q
        &\lesssim(sL)^{q/2}.
\end{split}
\end{equation}
\end{lemma}

\begin{proof}
For one row, the product structure in \eqref{eq:matched-score} gives
\begin{equation}\label{eq:row-score-factorization}
\begin{split}
        \norm{U^{(i)}}_s
        \lesssim{}&
        \left(
        \frac{|x_0^{(i)}|}{\sigma_0}
        +\frac{|r^{(i)}|}{\sigma_r}
        \right)
        \sup_{|J|\le s}
        \left\{(X_{-0,J}^{(i)})^\top
        C_{JJ}^{-1}X_{-0,J}^{(i)}\right\}^{1/2}\\
        &+\frac{\norm{d}_{C,s}}{\sigma_0}
        +\frac{\norm{h}_{C,s}}{\sigma_r}.
\end{split}
\end{equation}
For each $J$, a fixed net of the $C_{JJ}$-unit sphere has at most
$9^{|J|}$ points, and there are at most $(ep/s)^s$ supports.  The
relative sub-Gaussian assumption and integration of the resulting union
bound show that the supremum in \eqref{eq:row-score-factorization} has
$L_q$ norm $O(\sqrt{sL})$ for $q\le12$.  The two scalar factors have
bounded normalized $L_q$ norms, so H\"older's inequality and
\eqref{eq:proxy-energy} prove the first claim for $U^{(i)}$.  The
covariance of every normalized projection is uniformly bounded, so the
same support-net argument proves it for the matched Gaussian choice.

For the second claim, fix a net direction and let $V_i$ be its centered
scalar projection of $\widetilde U^{(i)}$.  This projection is either
Gaussian or a centered sum of products of sub-Gaussian variables, so
Bernstein's inequality gives
\begin{equation*}
        \Prob\left(
        \left|\frac1{\sqrt n}\sum_{i=1}^nV_i\right|>t
        \right)
        \le2\exp\{-c\min(t^2,t\sqrt n)\}
\end{equation*}
for each net direction, with a uniform constant.  A union over the same
supports and nets, followed by tail integration, gives an
$O(\sqrt{sL})$ bound in $L_q$ because $sL=o(n)$.  The proxy-energy
condition \eqref{eq:proxy-energy} makes the standardized sparse norm of
the deterministic mean $O(\sqrt{sL})$.
\end{proof}

\begin{theorem}[Gaussian approximation]
\label{thm:matched-chaos}
Suppose the regression model \eqref{eq:model} satisfies
\eqref{eq:proxy-energy} and
$sL=o(n)$.  Then
\begin{equation}\label{eq:empirical-chaos}
\begin{split}
        \sup_{S\subseteq[p],\,|S|=s}\widehat\beta_{0,S}
        &=\widehat\beta_{0,\emptyset}
        +\frac{T_C(Z_{1n},Z_{2n})}{n\sigma_0^2}
        +O_{\Prob}\!\left(
        \frac{\sigma_r}{\sigma_0}\delta_{s,n,p}^3
        \right).
\end{split}
\end{equation}

If in addition $(sL)^4=o(n)$, then the matched Gaussian vector in
\eqref{eq:matched-gaussian} satisfies
\begin{equation}\label{eq:smooth-approximation}
        d_3\!\left(
        \frac{n\sigma_0}{\sigma_rsL}
        \left\{
        \sup_{S\subseteq[p],\,|S|=s}\widehat\beta_{0,S}
        -\widehat\beta_{0,\emptyset}
        \right\},
        \frac{T_C(Z_1,Z_2)}{\sigma_0\sigma_rsL}
        \right)
        \lesssim
        \left(\frac{(sL)^4}{n}\right)^{1/6}
        +o(1).
\end{equation}
In particular, the two scalar laws can be coupled so that
\begin{equation}\label{eq:coupled-approximation}
        \sup_{S\subseteq[p],\,|S|=s}\widehat\beta_{0,S}
        =\widehat\beta_{0,\emptyset}
        +\frac{T_C(Z_1,Z_2)}{n\sigma_0^2}
        +o_{\Prob}\!\left(
        \frac{\sigma_r}{\sigma_0}\frac{sL}{n}
        \right).
\end{equation}
\end{theorem}

\begin{proof}
The empirical reduction follows from an exact regression identity.  With
$\widehat C=X_{-0}^\top X_{-0}/n$, block projection gives
\begin{equation}\label{eq:exact-regression}
        \widehat\beta_{0,S}-\widehat\beta_{0,\emptyset}
        =
        \frac{-Z_{1n,S}^\top\widehat C_{SS}^{-1}Z_{2n,S}
        +\dfrac{x_0^\top r}{x_0^\top x_0}
        Z_{1n,S}^\top\widehat C_{SS}^{-1}Z_{1n,S}}
        {x_0^\top x_0-Z_{1n,S}^\top
        \widehat C_{SS}^{-1}Z_{1n,S}}.
\end{equation}
The relative sub-Gaussian assumption and a sparse-net argument yield the
covariance-normalized block-Gram event
\begin{equation}\label{eq:block-gram}
        \sup_{|J|\le2s}
        \norm{C_{JJ}^{-1/2}\widehat C_{JJ}C_{JJ}^{-1/2}-I}_{\mathrm{op}}
        =O_{\Prob}(\delta_{s,n,p}).
\end{equation}
The same argument, now applied to the product scores, gives
\begin{equation}\label{eq:score-size}
        \norm{Z_{1n}}_{C,s}=O_{\Prob}(\sigma_0\sqrt{sL}),
        \qquad
        \norm{Z_{2n}}_{C,s}=O_{\Prob}(\sigma_r\sqrt{sL}).
\end{equation}
The proxy bound \eqref{eq:proxy-energy} covers the means in
\eqref{eq:score-size}.  The remaining scalar terms satisfy
\begin{equation*}
        x_0^\top x_0
        =n\sigma_0^2\{1+O_{\Prob}(n^{-1/2})\},
        \qquad
        \frac{x_0^\top r}{x_0^\top x_0}
        =O_{\Prob}\!\left(\frac{\sigma_r}{\sigma_0\sqrt n}\right),
\end{equation*}
because $x_0^{(i)}$ and $r^{(i)}$ are uncorrelated sub-Gaussian
variables.  Substitution into \eqref{eq:exact-regression}, followed by
the inverse perturbation bound on \eqref{eq:block-gram}, proves
\eqref{eq:empirical-chaos}.

For the Gaussian approximation, write
$\widetilde Q_S(z)=Q_S(z_1,z_2)/(\sigma_0\sigma_r)$ and smooth the
supremum with $\lambda>0$ by
\begin{equation*}
        f_\lambda(z)
        =\frac1\lambda
        \log\sum_{|S|=s}\exp\{\lambda\widetilde Q_S(z)\}.
\end{equation*}
Since $\log\binom ps\le sL$, replacing
$T_C(z_1,z_2)/(\sigma_0\sigma_rsL)$ by $f_\lambda(z)/(sL)$ costs at most
$1/\lambda$.

Each normalized support value obeys, for directions $u,v$,
\begin{equation*}
        |D\widetilde Q_S(z)[u]|\le \norm{z}_s\norm{u}_s,
        \qquad
        |D^2\widetilde Q_S(z)[u,v]|\le \norm{u}_s\norm{v}_s.
\end{equation*}
Softmax differentiation therefore gives
\begin{equation}\label{eq:softmax-derivatives}
\begin{aligned}
        |Df_\lambda(z)[u]|&\lesssim \norm{z}_s\norm{u}_s,\\
        |D^2f_\lambda(z)[u,u]|&\lesssim
        \{1+\lambda\norm{z}_s^2\}\norm{u}_s^2,\\
        |D^3f_\lambda(z)[u,u,u]|&\lesssim
        \{\lambda\norm{z}_s+\lambda^2\norm{z}_s^3\}\norm{u}_s^3.
\end{aligned}
\end{equation}
Let $\varphi$ belong to the class defining $d_3$ and apply Taylor's
formula to $\Phi(z)=\varphi\{f_\lambda(z)/(sL)\}$ while replacing the $n$
score rows by matched Gaussian rows one at a time.  The chain rule and
\eqref{eq:softmax-derivatives} give
\begin{equation}\label{eq:composite-third-derivative}
\begin{split}
        |D^3\Phi(z)[u,u,u]|
        \lesssim \norm{u}_s^3\bigg\{
        &\frac{\norm{z}_s^3}{(sL)^3}
        +\frac{\norm{z}_s}{(sL)^2}
        +\frac{\lambda\norm{z}_s^3}{(sL)^2}\\
        &+\frac{\lambda\norm{z}_s}{sL}
        +\frac{\lambda^2\norm{z}_s^3}{sL}
        \bigg\}.
\end{split}
\end{equation}
The first two Taylor terms cancel because the rows have the same mean and
covariance.  Lemma~\ref{lem:sparse-score-moments} and
\eqref{eq:composite-third-derivative} bound the sum of the third-order
remainders by
\begin{equation}\label{eq:lindeberg-rate}
        \frac{1+\lambda sL+\lambda^2(sL)^2}{\sqrt n}.
\end{equation}
Balancing \eqref{eq:lindeberg-rate} with the smoothing error using
$\lambda=\{\sqrt n/(sL)^2\}^{1/3}$ gives
$\{(sL)^4/n\}^{1/6}$.  This proves \eqref{eq:smooth-approximation} after
using \eqref{eq:empirical-chaos}.  The normalized statistics are tight
by \eqref{eq:score-size}; smoothing the scalar test functions and applying
the standard coupling characterization of weak convergence then yields
\eqref{eq:coupled-approximation}.
\end{proof}

A Gaussian multiplier estimates $\Sigma_U$ from the row scores.
Let $\bar U=n^{-1}\sum_iU^{(i)}$ and define
\begin{equation*}
        \widehat\Sigma_U
        =\frac1n\sum_{i=1}^n
        (U^{(i)}-\bar U)(U^{(i)}-\bar U)^\top.
\end{equation*}
Conditionally on the data, let $e_i\sim N(0,1)$ be independent and set
\begin{equation}\label{eq:multiplier}
        \widehat Z=\begin{pmatrix}\widehat Z_1\\\widehat Z_2\end{pmatrix}
        :=\begin{pmatrix}\sqrt n\,d\\\sqrt n\,h\end{pmatrix}
        +\frac1{\sqrt n}\sum_{i=1}^ne_i(U^{(i)}-\bar U).
\end{equation}
Conditionally on the observations, $\widehat Z$ is Gaussian with the same
mean as $Z$ and covariance $\widehat\Sigma_U$.  The construction is
oracle because $d$, $h$, and $r$ contain population quantities.

\begin{lemma}[Sparse multiplier covariance]
\label{lem:multiplier-covariance}
For $J\subseteq[p]$, let $U^{(i)}[J]$, $\Sigma_U[J]$, and
$\widehat\Sigma_U[J]$ contain both score channels indexed by $J$, and put
\begin{equation*}
        \mathcal W_J
        =\diag(\sigma_0^{-1}C_{JJ}^{-1/2},
        \sigma_r^{-1}C_{JJ}^{-1/2}).
\end{equation*}
Under the assumptions of Theorem~\ref{thm:matched-chaos},
\begin{equation}\label{eq:multiplier-covariance-rate}
\begin{split}
        \nu_n
        &:=
        \sup_{|J|\le2s}
        \norm{\mathcal W_J
        \{\widehat\Sigma_U[J]-\Sigma_U[J]\}
        \mathcal W_J^\top}_{\mathrm{op}}\\
        &=O_{\Prob}\!\left(
        \sqrt{\frac{sL}{n}}+\frac{(sL)^2}{n}
        \right).
\end{split}
\end{equation}
\end{lemma}

\begin{proof}
Every linear form of $\mathcal W_JU^{(i)}[J]$ is a sum of centered
products of normalized sub-Gaussian variables and therefore has uniformly
bounded $\psi_1$ norm.  Its centered square has bounded
$\psi_{1/2}$ norm.  A Fuk--Nagaev truncation bound for the average of
these squares, followed by a $1/4$-net on each sphere and a union over the
at most $\exp(c_1sL)$ support-net pairs, gives
\begin{equation*}
        \sqrt{\frac{sL+t}{n}}
        +\frac{(sL+t)^2}{n}
\end{equation*}
with probability at least $1-2e^{-t}$.  Fixed $t$ gives
\eqref{eq:multiplier-covariance-rate}; the sample-centering term
$\bar U\bar U^\top$ is smaller by the score-size bound.
\end{proof}

\begin{corollary}[Gaussian multiplier]
\label{cor:multiplier}
Under the assumptions of Theorem~\ref{thm:matched-chaos}, including
$(sL)^4=o(n)$, the conditional law of
\begin{equation}\label{eq:multiplier-statistic}
        \frac{T_C(\widehat Z_1,\widehat Z_2)}{\sigma_0\sigma_rsL}
\end{equation}
has $d_3$ distance $o_{\Prob}(1)$ from the law of the normalized searched
displacement in \eqref{eq:smooth-approximation}.
\end{corollary}

\begin{proof}
Lemma~\ref{lem:multiplier-covariance} gives $\nu_n=o_{\Prob}(1)$.
Gaussian interpolation between $\Sigma_U$ and $\widehat\Sigma_U$ uses the
smooth supremum from the proof of Theorem~\ref{thm:matched-chaos}.  With
$\alpha_S$ denoting the softmax weights, its Hessian is
\begin{equation*}
        \nabla^2 f_\lambda
        =\sum_S\alpha_S\nabla^2\widetilde Q_S
        +\lambda\operatorname{Cov}_\alpha(\nabla\widetilde Q_S).
\end{equation*}
Differentiating the interpolated expectation of
$\varphi\{f_\lambda/(sL)\}$ gives
\begin{equation}\label{eq:multiplier-interpolation}
\begin{split}
        \frac1{2sL}\E\left[
        \varphi'\operatorname{tr}\{(\widehat\Sigma_U-\Sigma_U)
        \nabla^2f_\lambda\}\right]
        +\frac1{2(sL)^2}\E\left[
        \varphi''(\nabla f_\lambda)^\top
        (\widehat\Sigma_U-\Sigma_U)\nabla f_\lambda\right].
\end{split}
\end{equation}
The supportwise Hessian is confined to one $s$-set, and every difference
of two softmax gradients is confined to the two score channels over the
union of two such sets.  The definition of $\nu_n$, the score-size bound,
and \eqref{eq:multiplier-interpolation} therefore bound the normalized
derivative by $O_{\Prob}(\nu_n+\nu_n\lambda)$.  The smoothing error is
$1/\lambda$, so $\lambda=\nu_n^{-1/2}$ gives a normalized $d_3$ error
$O_{\Prob}(\sqrt{\nu_n})$.  Combining this interpolation bound with
Theorem~\ref{thm:matched-chaos} proves the corollary.
\end{proof}

For jointly Gaussian rows, the score covariance is
\begin{equation}\label{eq:gaussian-score-covariance}
        \Sigma_U=
        \begin{pmatrix}
        \sigma_0^2C+dd^\top&hd^\top\\
        dh^\top&\sigma_r^2C+hh^\top
        \end{pmatrix}.
\end{equation}
Gaussian regression then couples the empirical and matched scores directly
under $sL=o(n)$.

\begin{corollary}[Gaussian rows]
\label{cor:gaussian}
If $(x_0^{(i)},r^{(i)},X_{-0}^{(i)})$ is jointly Gaussian, then under
\eqref{eq:proxy-energy} and $sL=o(n)$,
\begin{equation*}
        \sup_{S\subseteq[p],\,|S|=s}\widehat\beta_{0,S}
        =\widehat\beta_{0,\emptyset}
        +\frac{T_C(Z_1,Z_2)}{n\sigma_0^2}
        +O_{\Prob}\!\left(
        \frac{\sigma_r}{\sigma_0}\delta_{s,n,p}^3
        \right),
\end{equation*}
on an extension of the data probability space, with $\Sigma_U$ given by
\eqref{eq:gaussian-score-covariance}.
\end{corollary}

\begin{proof}
Gaussian regression writes the control row as a sum of its projections
onto $x_0^{(i)}$ and $r^{(i)}$ and an independent Gaussian remainder.
Conditional on the two scalar sample vectors, the remainder contributes a
matrix Gaussian whose $2\times2$ column covariance is their sample Gram
matrix.  Replacing that Gram matrix by the identity and coupling the three
scalar quadratic sums to Gaussian variables costs
$O_{\Prob}(\delta_{s,n,p})$ in the standardized sparse score norm.
The Lipschitz bound for \eqref{eq:search-functional}, together with
\eqref{eq:score-size}, makes the coupling error smaller than the block-Gram
error already present in \eqref{eq:empirical-chaos}.
\end{proof}

\begin{corollary}[Separable matched scores]
\label{cor:separable-matched-scores}
Suppose the assumptions of Theorem~\ref{thm:matched-chaos} hold, including
$(sL)^4=o(n)$ and $L\to\infty$, and the matched Gaussian scores satisfy
\eqref{eq:separable-gaussian-channels} with
$(a,b)=(Z_1,Z_2)$ and $(\sigma_a,\sigma_b)=(\sigma_0,\sigma_r)$.  If
\begin{equation}\label{eq:matched-passive-mean}
        n\norm{d/\sigma_0+h/\sigma_r}_{C,s}^2=o(sL),
\end{equation}
then the coupling in Theorem~\ref{thm:matched-chaos} satisfies
\begin{equation}\label{eq:matched-contrast-coupling}
        \sup_{|S|=s}\widehat\beta_{0,S}
        =\widehat\beta_{0,\emptyset}
        +\frac{\sigma_r}{2n\sigma_0}\norm{G}_{C,s}^2
        +o_{\Prob}\!\left(
        \frac{\sigma_r}{\sigma_0}\frac{sL}{n}
        \right).
\end{equation}
The corresponding distributional approximation is
\begin{equation}\label{eq:matched-contrast-distribution}
\begin{split}
        d_3\!\left(
        \frac{n\sigma_0}{\sigma_rsL}
        \left\{
        \sup_{|S|=s}\widehat\beta_{0,S}
        -\widehat\beta_{0,\emptyset}
        \right\},
        \frac{\norm{G}_{C,s}^2}{2sL}
        \right)=o(1).
\end{split}
\end{equation}
If a $G$-measurable support $\widehat S$ satisfies
\begin{equation}\label{eq:computed-projection-gap}
        \norm{G}_{C,s}^2
        -G_{\widehat S}^\top C_{\widehat S\widehat S}^{-1}
        G_{\widehat S}
        =o_{\Prob}(sL),
\end{equation}
then $Q_{\widehat S}(Z_1,Z_2)/(n\sigma_0^2)$ may replace the
contrast-energy term in \eqref{eq:matched-contrast-coupling}.
\end{corollary}

\begin{proof}
Under \eqref{eq:separable-gaussian-channels},
$\E H=\sqrt{n/2}\,(d/\sigma_0+h/\sigma_r)$.
Equations~\eqref{eq:contrast-oracle-bound},
\eqref{eq:passive-chi-square}, and
\eqref{eq:matched-passive-mean} make the contrast error
$o_{\Prob}(\sigma_0\sigma_rsL)$.  Substitution in
\eqref{eq:coupled-approximation} and
\eqref{eq:smooth-approximation} proves
\eqref{eq:matched-contrast-coupling} and
\eqref{eq:matched-contrast-distribution}.  Equation
\eqref{eq:contrast-transfer} proves the final claim.
\end{proof}

\begin{corollary}[Centered Gaussian rows]
\label{cor:centered-gaussian-rows}
Suppose $(x_0^{(i)},r^{(i)},X_{-0}^{(i)})$ is jointly Gaussian,
$d=h=0$, and $sL=o(n)$.  For every $C$ satisfying the standing sparse
block assumptions, let $G$ be the contrast in \eqref{eq:contrasts} formed
from $(Z_1,Z_2)$ with scales $(\sigma_0,\sigma_r)$.  No local-isotropy
condition is imposed, and on an extension of the data probability space,
\begin{equation}\label{eq:centered-gaussian-regression}
\begin{split}
        \sup_{|S|=s}\widehat\beta_{0,S}
        =\widehat\beta_{0,\emptyset}
        +\frac{\sigma_r}{2n\sigma_0}\norm{G}_{C,s}^2
        +O_{\Prob}\!\left[
        \frac{\sigma_r}{\sigma_0}
        \left\{\frac{s}{n}+\delta_{s,n,p}^3\right\}
        \right].
\end{split}
\end{equation}
\end{corollary}

\begin{proof}
When $d=h=0$, \eqref{eq:gaussian-score-covariance} is the Kronecker law
\eqref{eq:separable-gaussian-channels} with $\gamma=0$.
Corollary~\ref{cor:gaussian} reduces the regression search to $T_C$, and
\eqref{eq:centered-gaussian-contrast} replaces $T_C$ by the search
contrast energy with an $O_{\Prob}(\sigma_0\sigma_rs)$ remainder.
\end{proof}

Condition \eqref{eq:matched-passive-mean} is sharper than
\eqref{eq:proxy-energy}: the latter yields only an $O(sL)$ bound for the
remainder mean energy.  Corollary~\ref{cor:separable-matched-scores}
therefore requires a sharper little-$o$ bound, which may arise from
cancellation between the two standardized score means.
Corollary~\ref{cor:centered-gaussian-rows} has no such restriction because
both means vanish.

\section{Sparse projection}

The search contrast gives the sparse projection program
\begin{equation}\label{eq:sparse-projection-program}
        \norm{G}_{C,s}^2
        =\max_{\theta\in\R^p,\,\norm{\theta}_0\le s}
        \left\{2G^\top\theta-\theta^\top C\theta\right\}.
\end{equation}
An exact sparse-regression solver eliminates the projection gap in
\eqref{eq:contrast-transfer}.  More generally, if a solver returns a
support $\widehat S$ and an upper bound
$B\ge\norm{G}_{C,s}^2$, then
\begin{equation}\label{eq:computed-contrast-certificate}
\begin{split}
        0\le T_C(a,b)-Q_{\widehat S}(a,b)
        \le\frac{\sigma_a\sigma_b}{2}\bigg\{
        &B-G_{\widehat S}^\top C_{\widehat S\widehat S}^{-1}
        G_{\widehat S}\\
        &+H_{\widehat S}^\top C_{\widehat S\widehat S}^{-1}
        H_{\widehat S}\bigg\}.
\end{split}
\end{equation}
This separates the computed optimization error from the remainder energy.
For a finite family, braces with subscript $(t)$ below denote its $t$th
largest member.

A bounded restricted condition number gives a direct top-$s$
construction.  Write
\begin{equation}\label{eq:restricted-condition-number}
        \kappa_s(C)
        =\sup_{1\le|J|\le2s}
        \frac{\lambda_{\max}(C_{JJ})}{\lambda_{\min}(C_{JJ})},
\end{equation}
and let $S_{\mathrm{top}}$ index the largest $s$ values of $G_j^2$, with
ties resolved deterministically.

\begin{corollary}[Top-$s$ contrast]
\label{cor:top-s-contrast}
For every $C$,
\begin{equation}\label{eq:top-s-contrast-energy}
        G_{S_{\mathrm{top}}}^\top
        C_{S_{\mathrm{top}}S_{\mathrm{top}}}^{-1}G_{S_{\mathrm{top}}}
        \ge\frac{1}{\kappa_s(C)}\norm{G}_{C,s}^2.
\end{equation}
Consequently, the selected support value satisfies
\begin{equation}\label{eq:top-s-contrast-bound}
        \frac{T_C(a,b)}{\kappa_s(C)}
        -\frac{\sigma_a\sigma_b}{2}
        H_{S_{\mathrm{top}}}^\top
        C_{S_{\mathrm{top}}S_{\mathrm{top}}}^{-1}
        H_{S_{\mathrm{top}}}
        \le Q_{S_{\mathrm{top}}}(a,b)\le T_C(a,b).
\end{equation}
Under \eqref{eq:conditional-passive-covariance} with $K=O(1)$, the
remainder in \eqref{eq:top-s-contrast-bound} is
$O_{\Prob}\{\sigma_a\sigma_b
(s+\norm{\E[H\mid G]}_{C,s}^2)\}$.
\end{corollary}

\begin{proof}
Let $S_G$ maximize the projection energy and put
$J=S_{\mathrm{top}}\cup S_G$.  The definition of $S_{\mathrm{top}}$
gives $\norm{G_{S_{\mathrm{top}}}}_2^2\ge
\norm{G_{S_G}}_2^2$.  Interlacing then gives
\begin{equation*}
\begin{split}
        G_{S_{\mathrm{top}}}^\top
        C_{S_{\mathrm{top}}S_{\mathrm{top}}}^{-1}G_{S_{\mathrm{top}}}
        &\ge\frac{\norm{G_{S_{\mathrm{top}}}}_2^2}
        {\lambda_{\max}(C_{JJ})}\\
        &\ge\frac{\lambda_{\min}(C_{JJ})}
        {\lambda_{\max}(C_{JJ})}
        G_{S_G}^\top C_{S_GS_G}^{-1}G_{S_G},
\end{split}
\end{equation*}
which proves \eqref{eq:top-s-contrast-energy}.
The upper bound in Theorem~\ref{thm:contrast-approximation} and the
identity \eqref{eq:contrast-identity} prove
\eqref{eq:top-s-contrast-bound};
\eqref{eq:passive-contrast-size} gives its stochastic remainder.
\end{proof}

A covariance-aware search uses the exact marginal projection gains.  For
$v\in\R^p$, $S\subseteq[p]$ with $|S|<2s$, and $j\notin S$, define
\begin{equation}\label{eq:residual-score}
        v_{j\mid S}
        :=\frac{v_j-C_{jS}C_{SS}^{-1}v_S}
        {\{1-C_{jS}C_{SS}^{-1}C_{Sj}\}^{1/2}},
        \qquad v_{j\mid\emptyset}:=v_j.
\end{equation}
Block inversion gives
\begin{equation}\label{eq:projection-increment}
\begin{split}
        G_{S\cup\{j\}}^\top C_{S\cup\{j\},S\cup\{j\}}^{-1}
        G_{S\cup\{j\}}
        -G_S^\top C_{SS}^{-1}G_S
        =G_{j\mid S}^2.
\end{split}
\end{equation}
Forward projection therefore takes the largest squared residualized
coordinate at each step.  These coordinates also give an upper
certificate for any computed support.

\begin{corollary}[Projection certificate]
\label{cor:projection-certificate}
Suppose $2s\le p$.  For every $S$ with $|S|\le s$,
\begin{equation}\label{eq:projection-gap-certificate}
\begin{split}
        0\le\norm{G}_{C,s}^2-G_S^\top C_{SS}^{-1}G_S
        \le \kappa_s(C)\sum_{t=1}^s
        \{G_{j\mid S}^2:j\notin S\}_{(t)}.
\end{split}
\end{equation}
For a size-$s$ support, the right side of
\eqref{eq:contrast-transfer} is therefore at most
\begin{equation}\label{eq:computed-residual-certificate}
        \frac{\sigma_a\sigma_b}{2}
        \left\{
        \kappa_s(C)\sum_{t=1}^s
        \{G_{j\mid S}^2:j\notin S\}_{(t)}
        +H_S^\top C_{SS}^{-1}H_S
        \right\}.
\end{equation}
\end{corollary}

\begin{proof}
Let $A$ maximize the projection energy and put $R=A\setminus S$.  With
\begin{equation*}
        w=G_R-C_{RS}C_{SS}^{-1}G_S,
        \qquad
        \Gamma=C_{RR}-C_{RS}C_{SS}^{-1}C_{SR},
        \qquad D=\diag(\Gamma),
\end{equation*}
block inversion gives the joint gain $w^\top\Gamma^{-1}w$, while the sum
of the singleton gains is $w^\top D^{-1}w$.  The eigenvalues of the Schur
complement $\Gamma$ lie between the extreme eigenvalues of
$C_{A\cup S,A\cup S}$, and $D\preceq\lambda_{\max}(\Gamma)I$.  Therefore
\begin{equation*}
        w^\top D^{-1}w
        \ge\frac{1}{\kappa_s(C)}w^\top\Gamma^{-1}w.
\end{equation*}
Projection energy is monotone, so the joint gain is at least the gap
between the energies on $A$ and $S$.  The sum of the $s$ largest residual
gains dominates the singleton gains over $R$, proving
\eqref{eq:projection-gap-certificate}.  Combining this bound with
\eqref{eq:contrast-transfer} proves \eqref{eq:computed-residual-certificate}.
\end{proof}

Suppose $2s\le p$.  The forward-projection path is defined by
$J_0=\emptyset$ and
\begin{equation*}
        j_t\in\arg\max_{j\notin J_{t-1}}G_{j\mid J_{t-1}}^2,
        \qquad J_t=J_{t-1}\cup\{j_t\},
        \qquad 1\le t\le s,
\end{equation*}
with ties resolved deterministically.  Each quantity
\begin{equation}\label{eq:forward-upper-bound}
        G_{J_t}^\top C_{J_tJ_t}^{-1}G_{J_t}
        +\kappa_s(C)\sum_{\ell=1}^s
        \{G_{j\mid J_t}^2:j\notin J_t\}_{(\ell)}
        \qquad(0\le t\le s)
\end{equation}
is an upper bound for $\norm{G}_{C,s}^2$.  Their minimum can therefore be
used as $B$ in \eqref{eq:computed-contrast-certificate}.  When $C=I$,
the bound at $t=0$ is the top-$s$ square sum and equals the energy reached
by forward projection, so the computed optimization gap vanishes.  More
generally, whenever this certified gap is $o_{\Prob}(sL)$,
Corollary~\ref{cor:stochastic-contrast} makes
$Q_{J_s}(a,b)$ a stochastic approximation to $T_C(a,b)$.

\section{Product approximation}

The contrast approximation leaves the support geometry intact.  Under
local isotropy this geometry disappears and the sharper product formula
from the isotropic model is recovered.  When $C=I$,
\begin{equation}\label{eq:top-s-products}
        T_I(a,b)
        =\max_{|S|=s}\sum_{j\in S}(-a_jb_j)
        =\sum_{t=1}^s\{-a_jb_j:j\in[p]\}_{(t)}.
\end{equation}
This is the top-$s$ product form.  Under the unit-variance normalization,
local isotropy measures how far the reachable covariance blocks depart
from $I$:
\begin{equation}\label{eq:local-isotropy}
        \rho_s(C)
        :=\sup_{|J|\le2s}\norm{C_{JJ}-I}_{\mathrm{op}}.
\end{equation}

\begin{theorem}[Top-$s$ product approximation]
\label{thm:top-s-products}
For every $a,b\in\R^p$,
\begin{equation}\label{eq:static-product-bound}
        |T_C(a,b)-T_I(a,b)|
        \le\rho_s(C)\norm{a}_{C,s}\norm{b}_{C,s}.
\end{equation}
\end{theorem}

\begin{proof}
Fix $S$ with $|S|=s$ and put
$u=C_{SS}^{-1/2}a_S$ and $v=C_{SS}^{-1/2}b_S$.  Then
\begin{equation*}
        Q_S(a,b)+a_S^\top b_S
        =u^\top(C_{SS}-I)v,
\end{equation*}
whose absolute value is at most
$\rho_s(C)\norm{a}_{C,s}\norm{b}_{C,s}$.  Taking the uniform bound over
$S$ and using
$|\max_S x_S-\max_S y_S|\le\max_S|x_S-y_S|$ proves the result.
\end{proof}

\begin{corollary}[Matched Gaussian products]
\label{cor:matched-products}
Suppose the assumptions of Theorem~\ref{thm:matched-chaos} hold, including
$(sL)^4=o(n)$, and $\rho_s(C)=o(1)$.  The coupling in
Theorem~\ref{thm:matched-chaos} satisfies
\begin{equation}\label{eq:matched-product-coupling}
        \sup_{|S|=s}\widehat\beta_{0,S}
        =\widehat\beta_{0,\emptyset}
        +\frac{1}{n\sigma_0^2}
        \sum_{t=1}^s\{-Z_{1j}Z_{2j}:j\in[p]\}_{(t)}
        +o_{\Prob}\!\left(
        \frac{\sigma_r}{\sigma_0}\frac{sL}{n}
        \right).
\end{equation}
An independent matched Gaussian draw also gives
\begin{equation}\label{eq:matched-product-distribution}
\begin{split}
        d_3\!\left(
        \frac{n\sigma_0}{\sigma_rsL}
        \left\{
        \sup_{|S|=s}\widehat\beta_{0,S}
        -\widehat\beta_{0,\emptyset}
        \right\},
        \frac{1}{\sigma_0\sigma_rsL}
        \sum_{t=1}^s\{-Z_{1j}Z_{2j}:j\in[p]\}_{(t)}
        \right)=o(1).
\end{split}
\end{equation}
For jointly Gaussian rows, \eqref{eq:matched-product-coupling} and
\eqref{eq:matched-product-distribution} hold under the assumptions of
Corollary~\ref{cor:gaussian}, with $\rho_s(C)=o(1)$ and $sL=o(n)$.
\end{corollary}

\begin{proof}
Theorem~\ref{thm:top-s-products} and
Lemma~\ref{lem:sparse-score-moments} give
\begin{equation*}
        |T_C(Z_1,Z_2)-T_I(Z_1,Z_2)|
        =O_{\Prob}\{\rho_s(C)\sigma_0\sigma_rsL\}
        =o_{\Prob}(\sigma_0\sigma_rsL).
\end{equation*}
Substitution in \eqref{eq:coupled-approximation} and
\eqref{eq:smooth-approximation} proves
\eqref{eq:matched-product-coupling} and
\eqref{eq:matched-product-distribution}.  For Gaussian
rows, use Corollary~\ref{cor:gaussian} and
$\delta_{s,n,p}^3=o(sL/n)$.
\end{proof}

The second argument in \eqref{eq:matched-product-distribution} is the
top-$s$ sum of $p$ covariance-matched score products, which can be sorted
directly.  The products need not be independent because $\Sigma_U$ retains
the fourth-moment dependence of the data rows.  The covariance $C$
governs separability of the search, while $\Sigma_U$ governs dependence
among the product coordinates.  If
$\Sigma_U=\diag(\sigma_0^2I,\sigma_r^2I)$, all $2p$ Gaussian score
coordinates are mutually independent and the top-$s$ sum consists of
independent Gaussian products.

For the empirical scores, Theorem~\ref{thm:top-s-products} and
\eqref{eq:score-size} replace $T_C(Z_{1n},Z_{2n})$ in
\eqref{eq:empirical-chaos} by $T_I(Z_{1n},Z_{2n})$ whenever
$\rho_s(C)=o(1)$; \eqref{eq:top-s-products} identifies this value with the
empirical top-$s$ product sum.  For the multiplier scores,
Lemma~\ref{lem:multiplier-covariance} and a conditional Gaussian net bound
give the corresponding sparse-score bounds.  When $\rho_s(C)=o(1)$, the
same theorem replaces $T_C(\widehat Z_1,\widehat Z_2)$ in
\eqref{eq:multiplier-statistic} by
$T_I(\widehat Z_1,\widehat Z_2)$.
Bounded Lipschitz test functions transfer this conditional approximation
to the multiplier law.

\section{Relation to the design assumptions}

The whitened block-Gram event \eqref{eq:block-gram} is a relative
restricted-isometry property over $2s$-coordinate blocks, with error scale
$\delta_{s,n,p}$, and controls the reduction from the regression search to
$T_C$.  Replacing $Z_n$ by $Z$ requires $(sL)^4=o(n)$, while the direct
coupling for Gaussian rows requires only $sL=o(n)$.  This relative
restricted-isometry property is centered at $C$ and imposes no isotropy on
the population covariance.

Theorem~\ref{thm:contrast-approximation} is pointwise in $C$ and in the
two scores, so the relative restricted-isometry property enters only
through the regression reduction.
The stochastic remainder is governed instead by the conditional
covariance and mean of $H$.

The restricted condition number \eqref{eq:restricted-condition-number}
gives the sorting approximation in Corollary~\ref{cor:top-s-contrast} and
the certificate in Corollary~\ref{cor:projection-certificate} without
requiring $C$ to approach $I$.  Local isotropy is the stronger condition
$\rho_s(C)=o(1)$, under which
Theorem~\ref{thm:top-s-products} gives the top-$s$ product formula to
first order.

Under $\rho_s(C)=o(1)$, coordinatewise and invariant proxy bounds are
related by
\begin{equation}\label{eq:coordinate-to-invariant}
        \norm{v}_{C,2s}
        \le\frac{\sqrt{2s}\norm{v}_\infty}
        {\sqrt{1-\rho_s(C)}}
        \qquad(v\in\R^p).
\end{equation}
Write the residual outcome scale as
\begin{equation*}
        \varsigma^2
        =\Var(y^{(i)}-\beta_0x_0^{(i)})
        =\beta_{-0}^\top C\beta_{-0}+\Var(\varepsilon^{(i)}).
\end{equation*}
In particular, the coordinatewise conditions
\begin{equation*}
        \norm{d}_\infty\lesssim\sigma_0\sqrt{L/n},
        \qquad
        \norm{C\beta_{-0}}_\infty
        \lesssim\varsigma\sqrt{L/n}
\end{equation*}
imply
$\norm{d}_{C,2s}\lesssim\sigma_0\delta_{s,n,p}$ and
$\norm{C\beta_{-0}}_{C,2s}
\lesssim\varsigma\delta_{s,n,p}$ by
\eqref{eq:coordinate-to-invariant}.  The covariance-normalized conditions
also permit individual coordinates to be larger.
A partition of each $2s$-set into two $s$-sets gives
\begin{equation*}
        \norm{C\beta_{-0}}_{C,2s}^2
        \le\frac{2\{1+\rho_s(C)\}}{1-\rho_s(C)}
        \norm{C\beta_{-0}}_{C,s}^2.
\end{equation*}
Thus
$\norm{C\beta_{-0}}_{C,s}
\lesssim\varsigma\delta_{s,n,p}$ also yields
$\norm{C\beta_{-0}}_{C,2s}
\lesssim\varsigma\delta_{s,n,p}$.  If, in addition,
$|d^\top\beta_{-0}|=o(\sigma_0\varsigma\delta_{s,n,p})$, then
\eqref{eq:residual-outcome}, $sL=o(n)$, and the first bound in
\eqref{eq:proxy-energy} give
\begin{equation*}
        \frac{\sigma_r}{\varsigma}\longrightarrow1,
        \qquad
        \norm{h-C\beta_{-0}}_{C,2s}
        =o(\varsigma\delta_{s,n,p}^2).
\end{equation*}
Thus either
$\norm{C\beta_{-0}}_\infty
\lesssim\varsigma\sqrt{L/n}$ or
$\norm{C\beta_{-0}}_{C,s}
\lesssim\varsigma\delta_{s,n,p}$, together with
$|d^\top\beta_{-0}|
=o(\sigma_0\varsigma\delta_{s,n,p})$ and the bound on
$\norm{d}_{C,2s}$ in \eqref{eq:proxy-energy}, implies the corresponding
bound on $\norm{h}_{C,2s}$.

The Kronecker separation in \eqref{eq:separable-gaussian-channels} is not
a consequence of the relative restricted-isometry property.  The
deterministic contrast transfer remains valid without separation, while
\eqref{eq:conditional-passive-covariance} and
\eqref{eq:passive-mean-condition}, together with $L\to\infty$, suffice
for a negligible stochastic remainder.

\end{document}
